\begin{table}[H]
\centering
\caption{Top US Manufacturing Sectors by Gross Output: Liberation Day Tariff Context}
\label{tab:manufacturing_naics}
\begin{tabular}{p{4.5cm}rrl}
\toprule
Sector (NAICS) & GO 2021 & Share & Tariff Context \\
& (\$M) & of Mfg. & \\
\midrule
Petroleum refineries & 557,813 & 27.1\% & Energy/primary (10% floor) \\
Light truck and utility vehicle manufact.. & 367,006 & 17.8\% & Mixed: Canada/Mexico USMCA 10%, others higher \\
Pharmaceutical preparation manufacturing & 235,075 & 11.4\% & EU-weighted: ~20% avg tariff \\
Animal (except poultry) slaughtering, re.. & 169,319 & 8.2\% & Varies \\
Iron and steel mills and ferroalloy manu.. & 110,526 & 5.4\% & EU 20%, Canada 10%, Korea 25% \\
Aircraft manufacturing & 106,106 & 5.2\% & EU/Japan: ~20-24% \\
Other plastics product manufacturing & 104,942 & 5.1\% & China/Vietnam heavy exposure \\
Other basic organic chemical manufacturi.. & 103,181 & 5.0\% & EU and China exposure \\
Plastics material and resin manufacturin.. & 83,587 & 4.1\% & EU and China exposure \\
Printing & 75,634 & 3.7\% & Primarily domestic \\
Poultry processing & 72,883 & 3.5\% & Varies \\
Paperboard container manufacturing & 70,770 & 3.4\% & Varies \\
\midrule
\multicolumn{4}{l}{\textit{BLS Price Index Changes (2012--2018):}} \\
\bottomrule
\end{tabular}
\begin{tablenotes}
\small
\item Gross output from BEA via D\_GO\_by\_NAICS.xlsx (301 model data). Price indices from D\_price\_indices.xlsx (BLS PPI/MPI). Tariff context based on typical country-of-origin for each sector.
\end{tablenotes}
\end{table}